\section{Collar firmware}
\label{sec:firmware-collar}

The collar runs on Zephyr~\cite{zephyr}, a real-time operating system for microcontrollers, via the
nRF Connect SDK (NCS v3.3.1), targeting board
\texttt{xiao\_ble/nrf52840/sense}. It is a connectable BLE peripheral that senses on the cat,
advertises a compact telemetry packet, and streams paired audio\,+\,IMU clips for training once a
station connects and subscribes. It never touches Wi-Fi or the internet, so it stays small, low-power,
and battery-only. \texttt{build.ps1} drives the build: it locates the local NCS toolchain, runs
\texttt{west build -b xiao\_ble}, and copies out the UF2.

\subsection{Responsibilities}

The firmware has four jobs, driven from the main loop and BLE callbacks:
\begin{enumerate}[leftmargin=1.6em,itemsep=2pt]
  \item Run the rules-based activity gate: while the cat is active, sample the IMU and run the
        classifier cascade; on sustained stillness drop to low-power rest until motion returns.
  \item Advertise the result as an 8-byte BLE telemetry packet for the station to read
        (\cref{sec:protocols}).
  \item In training mode, stream paired audio\,+\,IMU frames to a subscribed station for capture.
  \item Receive new models over the air and swap them in safely.
\end{enumerate}

\subsection{Source layout}

\Cref{tab:collar-src} lists the collar firmware modules and what each is responsible for.

\begin{table}[ht]
  \centering
  \caption{Collar firmware modules (\texttt{firmware/collar/src/}).}
  \label{tab:collar-src}
  \begin{tabularx}{\textwidth}{@{}lX@{}}
    \toprule
    \textbf{File} & \textbf{Responsibility} \\
    \midrule
    \texttt{main.c}            & Boot and the telemetry loop; orchestration only \\
    \texttt{ble.c}             & Advertising, identity, the audio GATT service, notifications \\
    \texttt{audio.c}           & PDM microphone capture, decimated to 8\,kHz $\mu$-law \\
    \texttt{audio\_codec.h}    & $\mu$-law codec and the model-input audio representation \\
    \texttt{streaming.c}       & Decoupled reader/sender threads that assemble and send clip frames \\
    \texttt{imu.c}             & LSM6DS3TR-C continuous sampler \\
    \texttt{battery.c}         & 1S LiPo charge state via the onboard divider \\
    \texttt{classifier.c}      & The confidence-gated cascade policy (IMU first, audio confirms) \\
    \texttt{activity.c}        & The rules-based activity gate (cascade tier 0): still-timer that drives low-power rest \\
    \texttt{tflm\_classifier.cpp} & TensorFlow Lite Micro backend for the cascade \\
    \texttt{model\_loader.h}   & Runtime model-install API (\texttt{clf\_set\_imu\_model} / \texttt{clf\_set\_audio\_model}) \\
    \texttt{production.c}      & Rolling IMU window; runs inference and writes real telemetry \\
    \texttt{ota.c}             & OTA receive state machine; stages models to flash and defers the swap \\
    \bottomrule
  \end{tabularx}
\end{table}

The classifier is split deliberately. \texttt{classifier.c} defines the cascade policy in plain C
with \emph{weak} hooks for model-presence queries and inference calls; \texttt{tflm\_classifier.cpp}
provides the \emph{strong} \texttt{extern\,"C"} definitions that override them when the TensorFlow
Lite Micro backend is compiled in. This separates the policy from the heavy C++ inference, so the
firmware links and runs even with the backend or any model absent.

\subsection{Streaming and threading}

Capturing training data is the collar's most timing-sensitive path: PDM audio arrives steadily and
must be paired with IMU samples and pushed over BLE without dropping frames. The reader and sender
are therefore decoupled (\texttt{streaming.c}): one side assembles 100\,ms frames from the audio and
IMU sources, the other sends them over the audio characteristic, so a momentarily slow radio cannot
stall capture. The collar negotiates a large ATT MTU (247 bytes) and uses the 2\,Mbit PHY for
throughput.

\subsection{Activity gating and low-power rest}

The collar's power budget (NFR-1) is set by how much it senses, not by how fast it can classify, and a
cat is at rest for most of the day. A rules-based activity gate (\texttt{activity.c}), the first tier
of the cascade (\cref{sec:on-device-ai}), turns that idleness into battery life. Each production cycle
it takes a motion-energy proxy, the peak deviation of acceleration from 1\,g over the cycle's samples,
and runs a still-timer: once motion stays below a small threshold for a hold-off window (about a
minute), the collar enters \emph{low-power rest}.

In rest, the 104\,Hz sampling thread and all inference stop, the IMU drops to a low-power watch rate
(about 12.5\,Hz), and advertising slows to roughly a one-second interval; the main thread polls the
accelerometer slowly and the SoC idles between polls. A motion event above a wake threshold ends rest:
the IMU and advertising return to their active rates and the cascade resumes. Because the station logs
an episode whenever the collar's reported state changes, the collar advertises a dedicated
\texttt{REST} state byte while resting, so the entire dormant span is recorded as a single
\emph{rest} event covering exactly the period of low-to-no motion
(\cref{sec:protocols,sec:cloud-backend}). The thresholds and hold-off are constants in
\texttt{activity.h}, tuned against the measured power budget (\cref{sec:evaluation}).

\subsection{Model storage and flash layout}

Models are not linked into the firmware image. They live in a dedicated flash partition so a new model
can be written over the air without reflashing. The partition map is fixed in \texttt{pm\_static.yml}
to stop the partition manager from regrowing the application region over the model storage:

\begin{table}[ht]
  \centering
  \caption{Collar flash partitions (\texttt{pm\_static.yml}).}
  \label{tab:flash}
  \begin{tabularx}{\textwidth}{@{}llX@{}}
    \toprule
    \textbf{Partition} & \textbf{Address / size} & \textbf{Contents} \\
    \midrule
    \texttt{app}              & \texttt{0x27000} / \texttt{0x85000} & Application code (fixed size) \\
    \texttt{models\_storage}  & \texttt{0xAC000} / \texttt{0x40000} (256\,KiB) &
      Header plus the IMU (slot 0) and audio (slot 1) model blobs \\
    \texttt{settings\_storage} & \texttt{0xEC000} / \texttt{0x08000} & NVS settings \\
    \texttt{uf2\_partition}   & \texttt{0xF4000} / \texttt{0x0C000} & UF2 bootloader (unchanged) \\
    \bottomrule
  \end{tabularx}
\end{table}

\subsection{Build configuration}

On-device inference requires pulling TensorFlow Lite Micro into the NCS workspace (an optional,
normally-excluded west module) and turning on C++17 and the TFLM and CMSIS-NN options. Two stability
fixes are baked into the configuration: TFLM's \texttt{AllocateTensors()} and \texttt{Invoke()}
overflow Zephyr's small default stack, so the main stack is raised; and the model storage uses
\texttt{stream\_flash} to lazily erase pages during an OTA write.

\begin{lstlisting}[language=C,caption={Inference-related \texttt{prj.conf} settings (collar).},label={lst:prjconf}]
CONFIG_CPP=y
CONFIG_STD_CPP17=y
CONFIG_TENSORFLOW_LITE_MICRO=y
CONFIG_TENSORFLOW_LITE_MICRO_CMSIS_NN_KERNELS=y
CONFIG_MAIN_STACK_SIZE=16384      /* AllocateTensors()/Invoke() overflow the ~2 KiB default */
CONFIG_STREAM_FLASH_ERASE=y       /* lazily erase model pages during OTA */
CONFIG_BT_L2CAP_TX_MTU=247        /* large MTU for audio streaming + OTA chunks */
CONFIG_BT_CTLR_PHY_2M=y           /* 2 Mbit PHY for throughput */
\end{lstlisting}

The OTA receive path stages bytes to flash and defers the verify-and-swap onto the main thread rather
than committing a model from the BLE-RX callback (\cref{sec:protocols}).

\subsection{Development console and flashing}

The XIAO nRF52840 has no debug COM port, so flashing is always a UF2 drive-copy: double-tap RESET to
mount the bootloader drive, then re-run the build script to copy the UF2 across. For development the
firmware brings up a USB CDC ACM console on the \emph{legacy} USB device stack; the newer NEXT stack's
CDC console drops output on this board. The console prints a \texttt{MEOW> collar id=...} banner and a
periodic \texttt{state= steps= batt=} line; the dashboard reads the \texttt{id=} banner over Web
Serial to register the collar.

\subsection{Telemetry output}

In production the telemetry packet carries the cascade's live result: class and confidence from the
on-device classifier on the rolling IMU window, with the audio stage confirming uncertain results
(\cref{sec:on-device-ai}). It also reports an IMU-derived step count and the measured battery level.
While the activity gate holds the collar in low-power rest it advertises a dedicated \texttt{REST}
state instead, which the station records as a rest episode (\cref{sec:protocols}).
A freshly flashed collar advertises \texttt{UNKNOWN} until its first model arrives over the air, so
the firmware is always safe to run.
